\scanpage{183}
\mspage{495}
% position: top continued paragraph
\mbox{a u t r e} vie plus intense et plus délicate, que nous sommes trop endormis ou
trop pusillanimes souvent pour vivre dans la vie éveillée. Et il suffit de con-
templer un bébé endormi, ou seulement d'être tiré d'un sommeil profond, pour \unsure{sen|sen-}
tir que même sans rêves, le vrai bon sommeil est un \mbox{t r a v a i l} à sa fa-
çon\inserted{\unsure{:}} quelque chose qui nous absorbe de façon totale, à ``refaire le plein'' en
somme d'une énergie qui s'était dispersée et que nous venons \mbox{r e p u i s e r}
à sa source... C'est là, à nouveau, le ``yang dans le yin'', sans quoi le yin
lui-même serait destructeur.

% position: long middle paragraph
Des réflexions dans le même sens pourraient sûrement se développer aussi
pour l'inaction \mbox{é v e i l l é e}, en dehors du temps de sommeil. Il n'y a
qu'à observer sur pièces, de façon attentive, tel ou tel état qu'on perçoit
comme ``inaction''. On se rendra compte que dans l'inaction, il y a action, fût-
ce le stérile caquetage d'une pensée qui continue à tourner en rond alors \unsure{qu'ell|qu'elle}
a cessé de travailler. Mais à vrai dire, il est impropre d'appeler ``action'' ce
mouvement-là, purement mécanique, qui se poursuit par le seul effet d'une iner-
tie - par l'incapacité d'arrêter la machine! Et ce n'est certes pas cette agita-
tion intérieure qui va apporter à ``l'inaction'' une harmonie yin-yang qui la
rende bienfaisante. Par contre, il peut en être ainsi des activités diverses
destinées à meubler ses loisirs (quand ceux-ci sont néanmoins vécus comme un
état d'inaction). Mais même dans l'état de repos complet d'un état de convales-
cence disons, il peut y avoir action, sans quoi ce repos ou ``inaction'' devient
un \mbox{a v a c h i s s e m e n t}, peu propice certes à une convalescence (c'est
à-dire justement, au rétablissement d'un équilibre perturbé!). Par exemple,
cet état de repos peut susciter une attention à son propre corps et à son en-
tourage immédiat (qui en constitue comme une deuxième peau...), une prise de
connaissance donc voire une communion, qui par elle-même a un authentique ca-
ractère d'``action''; car nul doute que d' \mbox{a p p r e n d r e} est bel et bien
un \mbox{a c t e} (puisqu'il a un irrécusable \mbox{e f f e t}: l'apparition d'une
connaissance...).

% position: lower paragraph
En examinant un à un les quatorze couples que j'ai inclus dans le groupe
action-inaction (et sûrement on pourrait en trouver bien d'autres qui s'y insè-
rent naturellement), on constate que pour tous sauf peut-être un, c'est le pre-
mier terme, le terme ``masculin'', qui est investi de prestige, de ``valeur'', selon
les attitudes-réflexes véhiculés par notre culture et inculqués depuis l'enfance.
C'est le signe de toujours ce même déséquilibre invétéré dans notre culture,
% position: foot pencil
\margin{182}
